\creaCapitulo{Métodos de Resolución}
Para resolver los diferentes problemas hemos implementado tres métodos distintos, a saber: Runge-Kutta, que es explícito y utiliza pasos intermedios ($K1, K2$) para obtener los nuevos pasos; Adams-Bashford, que también es explícito y multipaso, pero utiliza pasos antiguos (como $f_{n-1}$ o $f_{n-2}$); y Adams-Bashford-Moulton, un predictor-corrector semiimplícito multipaso basado en Adams-Moulton (formalmente implícito) que utiliza Adams-Bashford para predecir $y_{n+1}$ y sustituirlo en $f(t_{n+1}, y_{n+1})$.


\creaSeccion{Runge-Kutta}

Ahora que hemos tratado el método de Euler, podremos dar el siguiente paso hacia Runge-Kutta. Aunque técnicamente ya lo hemos estado viendo, ya que los métodos de Runge-Kutta son generalizaciones de la fórmula básica de Euler. Por esta razón, se dice que el método de Euler es un método de Runge-Kutta de primer orden. \cite{Segarra2020}

En cada uno de los órdenes de Runge-Kutta resolveremos una serie de pasos intermedios para hallar el valor final. El de orden 1, como hemos mencionado, coincidiría con Euler, como se puede ver en la ecuación \ref{eq:EulerExplicito}.

Para ilustrar este punto, procedemos a pasar a la fórmula de segundo orden, en el que buscaremos encontrar constantes o parámetros $w_1, w_2, \alpha, \beta$ tal que la fórmula concuerda con un polinomio de Taylor de grado 2. \cite{Zill2009}
\begin{equation}
	\label{eq:Runge-Kutta2Inicial}
	\begin{split}
		y_{n+1} &= y_n + h \cdot (w_1 k_1 + w_2 k_2)\\
			k_1 &= f(x_n, y_n)\\
			k_2 &= f(x_n + \alpha h, y_n + \beta h k_1)
	\end{split}
\end{equation}
Para el uso que haremos en este proyecto, basta decir que esto se puede hacer siempre que las constantes cumplan que:
\begin{equation*}
	\begin{split}
		w_1 + w_2 = 1 \text{ , } w_2 \alpha = \frac{1}{2} \text{ y } w_2 \beta = \frac{1}{2} \\
		w_1 = 1 - w_2 \text{ , } \alpha = \frac{2}{2w_2} \text{ y } \beta = \frac{1}{2w_2}
	\end{split}
\end{equation*}
Este es un sistema algebraico que posee tres ecuaciones y cuatro incógnitas, lo cual implica un número infinito de soluciones, donde $w_2 \neq 0$, y para nuestra aplicación práctica tomaremos $w_1 = w_2 = \frac{1}{2}$ y $\alpha = \beta = 1$, formando por tanto:
\begin{equation}
	\label{eq:Runge-Kutta2}
	\begin{split}
		y_{n+1} &= y_n + h/2 \cdot (k_1 + k_2)\\
		k_1 &= f(t_n, y_n)\\
		k_2 &= f(t_n + h, y_n + k_1 \cdot h)
	\end{split}
\end{equation}

De forma análoga, obtendremos el tercer orden:
\begin{equation}
	\label{eq:Runge-Kutta3}
	\begin{split}
		y_{n+1} &= y_n + h/6 \cdot (k_1 +4k_2 + k_3)\\
		k_1 &= f(t_n, y_n)\\
		k_2 &= f(t_n + h/2, y_n + k_1 \cdot h/2) \\
		k_3 &= f(t_n + h, y_n - k_1 \cdot h + 2k_3 \cdot h)
	\end{split}
\end{equation}

En el método de Runge-Kutta de orden 4 resolveremos diversos pasos que nos servirán para estimar de forma muy aproximada el resultado. Será en éste en el que nos centraremos a la hora de las mediciones. Partiendo de un valor inicial $dy/dt = f(t,y)$, $y(t_0)=y_0$, y de un valor de salto $h>0$, calcularemos iterativamente los valores tal que:
\begin{equation}
	\label{eq:Runge-Kutta4}
	\begin{split}
		y_{n+1} &= y_n + h/6 \cdot (k_1 +2k_2 + 2k_3 + k_4)\\
		k_1 &= f(t_n, y_n)\\
		k_2 &= f(t_n + h/2, y_n + k_1 \cdot h/2) \\
		k_3 &= f(t_n + h/2, y_n + k_2 \cdot h/2) \\
		k_4 &= f(t_n + h, y_n + k_3 \cdot h)
	\end{split}
\end{equation}

Aplicado en un contexto de programación, tendremos una clase $RungeKutta$ con una función principal que aplicará el cálculo y algunas auxiliares que se explicarán más adelante. Dicha función principal tendrá el siguiente diseño, en el que se harán diversas llamadas a la evaluación de la derivada de cada problema en cuestión (aquí lo llamaremos \textit{feval}): 

\begin{algorithm}[H] %Le ponemos la H para que se coloque entre los dos párrafos, y no se mueva libremente
	\begin{algorithmic}
		\State $y_n \gets y_0$
		\State $t_n \gets t_0$
		\While{$t_n < t_f$}
		\State $K1 \gets feval(t_n, Y_n)$
		\State $K2 \gets feval(t_n+h/2, Y_n + K1\cdot h/2)$
		\State $K3 \gets feval(t_n+h/2, Y_n + K2\cdot h/2)$
		\State $K4 \gets feval(t_n+h, Y_n + K3\cdot h)$
		\State $Y_n \gets Y_n + h/6 \cdot (K1 + 2 \cdot K2 + 2 \cdot K3 + K4)$
		\State $t_n \gets t_n + h$
		\EndWhile
	\end{algorithmic}
	\caption{Runge-Kutta de Orden 4}
	\label{alg:RungeKutta4}
\end{algorithm}

Como podemos ver, la aplicación del algoritmo es relativamente fácil de entender, se trabaja con 4 vectores auxiliares en los que hacemos los cálculos intermedios, que luego se suman ponderadamente. Como cada evaluación K depende de la anterior, no podremos hacer una operación vectorial en la que se evalúen los cuatro K en paralelo, sino que será en las otras operaciones, las evaluaciones de cada miembro, las sumas, y las multiplicaciones de cada vector, en las que podremos entrar a trabajar el paralelismo para mejorar la eficiencia de nuestro código. 

\creaSeccion{Adams-Bashford}
A diferencia de Runge-Kutta, donde sólo considerábamos un valor $y_n$ para obtener $y_{n+1}$, en los métodos de Adams-Bashford, al ser multipaso, sí tendremos en cuenta más puntos, como $y_{n-1}$ o $y_{n-3}$. En este proyecto nos hemos centrado en los métodos de paso fijo, es decir, que la diferencia entre $y_n$ y $y_{n-1}$ es constante (e igual a $\Delta t$). Si bien existen otros métodos de paso variable, no serán tratados.

De forma general, tenemos que un método lineal de k pasos viene determinado por una expresión de la forma: \cite{Molero2007}
\begin{equation}
	\label{eq:MetodoLinealKPasos}
	\sum_{j=0}^{k} \alpha_j \cdot z_{n+j} = h \cdot \sum_{j=0}^{k} \beta_j \cdot f(x_{n+j} , z_{n+j}) , n\geq 0
\end{equation}

\creaSeccion{Adams-Bashford-Moulton}